% Options for packages loaded elsewhere
\PassOptionsToPackage{unicode}{hyperref}
\PassOptionsToPackage{hyphens}{url}
%
\documentclass[
  11pt,
]{article}
\usepackage{amsmath,amssymb}
\usepackage{iftex}
\ifPDFTeX
  \usepackage[T1]{fontenc}
  \usepackage[utf8]{inputenc}
  \usepackage{textcomp} % provide euro and other symbols
\else % if luatex or xetex
  \usepackage{unicode-math} % this also loads fontspec
  \defaultfontfeatures{Scale=MatchLowercase}
  \defaultfontfeatures[\rmfamily]{Ligatures=TeX,Scale=1}
\fi
\usepackage{lmodern}
\ifPDFTeX\else
  % xetex/luatex font selection
\fi
% Use upquote if available, for straight quotes in verbatim environments
\IfFileExists{upquote.sty}{\usepackage{upquote}}{}
\IfFileExists{microtype.sty}{% use microtype if available
  \usepackage[]{microtype}
  \UseMicrotypeSet[protrusion]{basicmath} % disable protrusion for tt fonts
}{}
\makeatletter
\@ifundefined{KOMAClassName}{% if non-KOMA class
  \IfFileExists{parskip.sty}{%
    \usepackage{parskip}
  }{% else
    \setlength{\parindent}{0pt}
    \setlength{\parskip}{6pt plus 2pt minus 1pt}}
}{% if KOMA class
  \KOMAoptions{parskip=half}}
\makeatother
\usepackage{xcolor}
\usepackage[margin=1in]{geometry}
\setlength{\emergencystretch}{3em} % prevent overfull lines
\providecommand{\tightlist}{%
  \setlength{\itemsep}{0pt}\setlength{\parskip}{0pt}}
\setcounter{secnumdepth}{-\maxdimen} % remove section numbering
\ifLuaTeX
  \usepackage{selnolig}  % disable illegal ligatures
\fi
\IfFileExists{bookmark.sty}{\usepackage{bookmark}}{\usepackage{hyperref}}
\IfFileExists{xurl.sty}{\usepackage{xurl}}{} % add URL line breaks if available
\urlstyle{same}
\hypersetup{
  pdftitle={A literal-window construction-quality exponent for Catalan's constant},
  pdfauthor={Working paper},
  hidelinks,
  pdfcreator={LaTeX via pandoc}}

\title{A literal-window construction-quality exponent for Catalan's
constant}
\author{Working paper}
\date{12 August 2026}

\begin{document}
\maketitle

\hypertarget{paper-theorem-and-proof-ledger-for-the-literal-windows}{%
\section{Paper theorem and proof ledger for the literal
windows}\label{paper-theorem-and-proof-ledger-for-the-literal-windows}}

\textbf{Status (12 August 2026).} This document gives a paper-level
proof, using the literal prime windows only, that an explicitly selected
sequence of integer combinations of the Zudilin and Nesterenko rows has
construction quality at least

\[
  0.581983473693357\ldots
\]

The new ingredient is the unconditional two-minimum selection theorem in
Section 6. It supplies a lower bound for the actual height and an upper
bound for the linear-form value relative to that height. No conditional
Lean theorem, finite certificate, numerical scan, or unproved common-gcd
estimate is used.

\hypertarget{quality-of-a-constructed-sequence}{%
\subsection{1. Quality of a constructed
sequence}\label{quality-of-a-constructed-sequence}}

Put

\[
 G=\sum_{k=0}^{\infty}\frac{(-1)^k}{(2k+1)^2}.
\]

An \textbf{admissible approximation sequence} to \(G\), indexed by an
infinite set \(I\subseteq\mathbb N\), is a sequence
\((q_n,p_n)\in\mathbb Z^2\), \(n\in I\), such that

\[
 q_n\ne0,\qquad q_nG-p_n\ne0,\qquad |q_n|\longrightarrow\infty.
\]

Its \textbf{construction quality} is

\[
 \delta_I(q,p)
 :=\liminf_{\substack{n\to\infty\\n\in I}}
 \frac{-\log|G-p_n/q_n|}{\log|q_n|}.
 \tag{1.1}
\]

When \(I\) is the set of all sufficiently large integers, we abbreviate
this quantity to \(\delta(q,p)\); changing finitely many terms has no
effect.

This is a property of the displayed family, not the classical
irrationality exponent of \(G\) (whose irrationality is unknown). If

\[
 \log|q_n|=An+o(n),\qquad
 \log|q_nG-p_n|=En+o(n),\qquad A>0,
 \tag{1.2}
\]

then

\[
 \delta_I(q,p)=1-\frac EA.
 \tag{1.3}
\]

More generally, a lower height estimate \(\log|q_n|\ge (A-o(1))n\),
together with the corresponding upper estimates, is indispensable when
\(E>0\). Upper bounds for both quantities do not imply (1.3). In this
normalization the irrationality threshold is \(\delta>1\), not merely
\(\delta>0\).

\hypertarget{the-two-rational-rows-and-explicit-integer-clearings}{%
\subsection{2. The two rational rows and explicit integer
clearings}\label{the-two-rational-rows-and-explicit-integer-clearings}}

Write \(D_m=\operatorname{lcm}(1,\ldots,m)\), and put
\(\varphi(t)=20t^2-8t+1\). Let \(Q_m,P_m\in\mathbb Q\) be the two
solutions of

\[
 \begin{split}
 &(2m+1)^2(2m+2)^2\varphi(m)z_{m+1}\\
 &\quad=(3520m^6+5632m^5+2064m^4-384m^3-156m^2+16m+7)z_m\\
 &\qquad +(2m-1)^2(2m)^2\varphi(m+1)z_{m-1}\qquad(m\ge1),
 \end{split}
 \tag{2.1}
\]

with

\[
 (Q_0,Q_1)=(1,7/4),\qquad(P_0,P_1)=(0,13/8).
 \tag{2.2}
\]

Set \(\Lambda_m^Z=Q_mG-P_m\). For \(m\ge1\), define the completely
explicit exponent

\[
 e_m=\min\{6m,\,4m+3+\lfloor\log_2(2m-1)\rfloor\}
 \tag{2.3}
\]

and, at \(m=3n\), define

\[
 X_n=2^{e_{3n}}D_{6n-1}^2Q_{3n},\qquad
 Y_n=2^{e_{3n}}D_{6n-1}^2P_{3n}.
 \tag{2.4}
\]

These are integers. Indeed, Zudilin's 2003 theorem gives

\[
 2^{4m+3}D_mQ_m\in\mathbb Z,qquad
 2^{4m+3}D_{2m-1}^3P_m\in\mathbb Z,
 \tag{2.5}
\]

whereas the partial-fraction argument in his second 2002 paper gives

\[
 2^{6m}Q_m\in\mathbb Z,qquad
 2^{6m}D_{2m-1}^2P_m\in\mathbb Z.
 \tag{2.6}
\]

If a rational number \(r\) is cleared both by \(2^a\) and by \(2^bM\),
then it is cleared by \(2^{\min(a,b+v_2(M))}\): the first hypothesis
says that the reduced denominator of \(r\) is a power of two, and the
second bounds its exponent by \(b+v_2(M)\). Apply this first to
\(r=D_{2m-1}^2P_m\), with \(a=6m\), \(b=4m+3\), and \(M=D_{2m-1}\). This
proves that \(2^{e_m}D_{2m-1}^2P_m\) is integral. For \(Q_m\), the same
argument with \(M=D_m\) gives the (possibly smaller) exponent \[
 \min\{6m,\,4m+3+\lfloor\log_2m\rfloor\}\le e_m;
\] multiplication by the remaining power of two proves that
\(2^{e_m}Q_m\) is integral. This gives (2.3)--(2.4). Equivalently,

\[
 2^{e_m}=2^{4m}s_m,qquad
 s_m=2^{\min\{2m,\,3+\lfloor\log_2(2m-1)\rfloor\}},
 \quad \log s_m=O(\log m).
 \tag{2.7}
\]

Thus the formerly implicit \(2^{o(m)}\) is the explicit factor \(s_m\).

For the Nesterenko row, let \((a)_j=a(a+1)\cdots(a+j-1)\). For \(n\ge1\)
define

\[
 \mathcal R_n(t)=
 \frac{(3n)!(3/2)_{3n-1}}{(4n)!}
 \frac{t(t-1)\cdots(t-4n+1)}
      {\prod_{r=0}^{3n}(t+r+1/2)^2}.
 \tag{2.8}
\]

The rational numbers \(A_{1,n,j},A_{2,n,j}\) are defined uniquely by the
partial-fraction identity

\[
 \mathcal R_n(t)=\sum_{j=0}^{3n}
 \frac{(t+1)\cdots(t+j)}{(t+1/2)_{j+1}}
 \left(A_{1,n,j}+\frac{A_{2,n,j}}{t+j+1/2}\right).
 \tag{2.9}
\]

In particular, taking the double-pole coefficient gives the explicit
positive formula

\[
 A_{2,n,j}=2^{-14n+2j+1}
 \frac{(8n+2j)!\,j!\,(6n)!}
 {(4n)!\,(4n+j)!\,(3n-j)!^2\,(2j)!^2}
 \quad(0\le j\le3n).
 \tag{2.9a}
\]

(An empty numerator product is \(1\).) Define

\[
 B_n=\sum_{j=0}^{3n}A_{2,n,j},
 \tag{2.10}
\]

\[
 C_n=\sum_{j=1}^{3n}\sum_{\nu=0}^{j-1}
 \frac{\nu!}{(3/2)_\nu}
 \left(A_{1,n,j}+\frac{A_{2,n,j}}{\nu+1/2}\right),
 \tag{2.11}
\]

and

\[
 J_n=\int_0^1\!\int_0^1
 \frac{x^{4n-1/2}(1-x)^{3n}y^{4n}(1-y)^{3n-1/2}}
 {(1-xy)^{4n+1}}\,dx\,dy.
 \tag{2.12}
\]

Nesterenko's partial-fraction theorem gives the exact identity

\[
 J_n=4B_nG-C_n.
 \tag{2.13}
\]

The correct integer row in this normalization is

\[
 V_n=4^{7n+1}D_{6n}^2B_n,qquad
 U_n=4^{7n}D_{6n}^2C_n,
 \tag{2.14}
\]

and

\[
 \Lambda_n^N=V_nG-U_n=4^{7n}D_{6n}^2J_n.
 \tag{2.15}
\]

Both \(V_n\) and \(U_n\) are integers. This is the conservative,
unconditional denominator theorem on page 171 of Nesterenko's paper; it
is not the later \(D_{8n}D_{6n}/(\Delta_n\Delta_{1,n})\) normalization
whose second coordinate is conjectural. The extra factor \(4\) in
\(V_n\) is forced by (2.13). It is missing from the displayed \(q_n\) in
the printed equation (4.4), but it reappears correctly in the paper's
(v\_n) on page 175.

\hypertarget{two-sided-asymptotics-in-these-normalizations}{%
\subsection{3. Two-sided asymptotics in these
normalizations}\label{two-sided-asymptotics-in-these-normalizations}}

Let

\[
 \phi=\frac{1+\sqrt5}{2},\qquad
 T=\frac{3303+437\sqrt{57}}{144},\qquad
 T^- =\frac{3303-437\sqrt{57}}{144}.
 \tag{3.1}
\]

Then

\[
 \begin{aligned}
 A_1&=12\log2+12+15\log\phi
       =27.535943542613395425\ldots,\\
 E_1&=12\log2+12-15\log\phi
       =13.099588790825292001\ldots,\\
 A_2&=14\log2+12+2\log T
       =29.354773862079159628\ldots,\\
 E_2&=14\log2+12+2\log T^-
       =14.393145267190103982\ldots.
 \end{aligned}
 \tag{3.2}
\]

In the exact normalizations (2.4) and (2.14), the primary-source limit
statements read, before clearing,

\[
\begin{array}{ll}
 n^{-1}\log Q_{3n}\to15\log\phi,&
 n^{-1}\log|Q_{3n}G-P_{3n}|\to-15\log\phi,\\
 n^{-1}\log B_n\to2\log T,&
 n^{-1}\log J_n\to2\log T^-.
\end{array}
\]

Moreover, (2.7) gives \(e_{3n}\log2=12n\log2+O(\log n)\), while
\(2\log D_{6n-1}=12n+o(n)\) and \(2\log D_{6n}=12n+o(n)\). Combining
these identities gives

\[
 \begin{array}{ll}
 \log X_n=A_1n+o(n),&\log Y_n=A_1n+o(n),\\
 \log V_n=A_2n+o(n),&\log U_n=A_2n+o(n),\\
 \log|\Lambda_n^Z|=E_1n+o(n),&
 \log|\Lambda_n^N|=E_2n+o(n).
 \end{array}
 \tag{3.3}
\]

Here all four coordinates are eventually positive. For the two second
coordinates this follows from \(P_m/Q_m\to G>0\) and
\(C_n/(4B_nG)\to1\), respectively; the latter limit follows from (2.13)
and \(T^-<T\). Thus (3.3) is two-sided, not just a collection of upper
estimates.

\hypertarget{literal-window-divisor-and-the-exact-congruence-lattice}{%
\subsection{4. Literal-window divisor and the exact congruence
lattice}\label{literal-window-divisor-and-the-exact-congruence-lattice}}

Define

\[
 S_n=\prod_{\substack{4n<\ell<6n\\ \ell\ \mathrm{prime}}}\ell^2
     \prod_{\substack{2n<\ell<3n\\ \ell\ \mathrm{prime}}}\ell^2.
 \tag{4.1}
\]

The two intervals are disjoint, and every prime in them occurs in both
\(D_{6n-1}\) and \(D_{6n}\). Hence

\[
 S_n\mid X_n,qquad S_n\mid V_n.
 \tag{4.2}
\]

The prime number theorem in the form \(\vartheta(x)=x+o(x)\) gives

\[
 \log S_n
 =2(\vartheta(6n)-\vartheta(4n))
  +2(\vartheta(3n)-\vartheta(2n))=6n+o(n).
 \tag{4.3}
\]

Put

\[
 L_n=\{(c_1,c_2)\in\mathbb Z^2:
        c_1Y_n+c_2U_n\equiv0\pmod{S_n}\}.
 \tag{4.4}
\]

Let

\[
 g_n=\gcd(S_n,Y_n,U_n),\qquad d_n=S_n/g_n.
 \tag{4.5}
\]

The homomorphism

\[
 (c_1,c_2)\longmapsto c_1Y_n+c_2U_n\pmod{S_n}
\]

has image the subgroup generated by \(g_n\pmod {S_n}\), of order
\(S_n/g_n\). Therefore

\[
 [\mathbb Z^2:L_n]=d_n=\frac{S_n}{\gcd(S_n,Y_n,U_n)}\le S_n.
 \tag{4.6}
\]

For every \(c=(c_1,c_2)\in L_n\),

\[
 q_n(c)=\frac{c_1X_n+c_2V_n}{S_n},\qquad
 p_n(c)=\frac{c_1Y_n+c_2U_n}{S_n}
 \tag{4.7}
\]

are integers and

\[
 q_n(c)G-p_n(c)=
 \frac{c_1\Lambda_n^Z+c_2\Lambda_n^N}{S_n}.
 \tag{4.8}
\]

\begin{quote}
\textbf{Main theorem.} Let \(G\), \(D_m\), \(Q_m\), \(P_m\), \(e_m\),
\(X_n\), \(Y_n\), \(\mathcal R_n\), \(A_{1,n,j}\), \(A_{2,n,j}\),
\(B_n\), \(C_n\), \(J_n\), \(V_n\), \(U_n\), \(S_n\), \(L_n\),
\(q_n(c)\), and \(p_n(c)\) be exactly the sequences defined in
(1.1)--(4.8). For every sufficiently large \(n\) there is a vector
\(c_n\in L_n\) such that the integer sequence \((q_n(c_n),p_n(c_n))\) is
admissible and has construction quality

\[
\delta(q,p)\ge0.5819834736933570952\ldots .          \tag{4.9}
\]

More exactly, the right side is

\[
1-\frac{E_{\rm eff}}{A_{\rm eff}},\qquad
A_{\rm eff}=A_2-\frac{6+E_2-E_1}{2},\quad
E_{\rm eff}=E_2-\frac{6+E_2-E_1}{2},               \tag{4.10}
\]

with \(A_i,E_i\) as in (3.2). This is a construction-quality statement;
it is not a claim that \(G\) is irrational.
\end{quote}

The rest of the paper proves this theorem. The selection of \(c_n\) is
given in Theorem 6.2 and includes the required lower height bound.

\hypertarget{eventual-independence-from-crossed-rates}{%
\subsection{5. Eventual independence from crossed
rates}\label{eventual-independence-from-crossed-rates}}

Let

\[
 \mathcal D_n=X_nU_n-V_nY_n.
 \tag{5.1}
\]

Using the two linear forms,

\[
 \mathcal D_n=V_n\Lambda_n^Z-X_n\Lambda_n^N.
 \tag{5.2}
\]

The two terms have rates \(A_2+E_1\) and \(A_1+E_2\), and

\[
 \gamma=(A_2+E_1)-(A_1+E_2)
 =0.52527384310095222099\ldots>0.
 \tag{5.3}
\]

It follows from the reverse triangle inequality and (3.3) that

\[
 \log|\mathcal D_n|=(A_2+E_1)n+o(n),
 \tag{5.4}
\]

so \(\mathcal D_n\ne0\) for all sufficiently large \(n\). This proves
eventual row independence; no finite determinant scan is being used.

\hypertarget{selection-theorem-and-construction-quality-bound}{%
\subsection{6. Selection theorem and construction-quality
bound}\label{selection-theorem-and-construction-quality-bound}}

Put \(K=6\) and

\[
 x_* =\frac{K+E_2-E_1}{2}
      =3.6467782381824059906\ldots.
 \tag{6.1}
\]

For an integer \(m\), write \(\log^+|m|=\log\max\{1,|m|\}\).

The following annular statement was the first sufficient formulation of
the selection problem. It is recorded to explain the rate ledger, but it
is not an assumption of the final theorem.

\begin{quote}
\textbf{Annular selection statement (AS).} There are numbers
\(x_n<x_*\), with \(x_n\to x_*\), and vectors
\(c_n=(c_{1,n},c_{2,n})\in L_n\), such that \[
\log^+|c_{1,n}|\le x_nn+o(n),\qquad
\log|c_{2,n}|=(K-x_n)n+o(n),                 \tag{6.2}
\] where the remainders in (3.3), (4.3), and (6.2) are
\(o((x_*-x_n)n)\). In particular \(c_{2,n}\ne0\).
\end{quote}

The separation condition is a standard diagonal formulation: it merely
says that \(x_n\uparrow x_*\) slowly enough for exponential dominance to
remain visible.

\begin{quote}
\textbf{Reduction theorem.} If (AS) holds for the explicitly defined
lattices (4.4), then the pairs \((q_n(c_n),p_n(c_n))\) are eventually
admissible and have \[
\begin{aligned}
\log|q_n(c_n)|&=A_{\rm eff}n+o(n),\\
\log|q_n(c_n)G-p_n(c_n)|&=E_{\rm eff}n+o(n),
\end{aligned}                                      \tag{6.3}
\] where \[
\begin{aligned}
A_{\rm eff}&=A_2-x_*=25.707995623896753637\ldots,\\
E_{\rm eff}&=E_2-x_*=10.746367029007697991\ldots.
\end{aligned}                                      \tag{6.4}
\] Consequently their construction quality is \[
\boxed{\delta_*=1-\frac{E_{\rm eff}}{A_{\rm eff}}
=0.5819834736933570952\ldots.}                     \tag{6.5}
\] In fact the quotient in (1.1) has a limit, so its liminf and limsup
are both equal to \(\delta_*\).
\end{quote}

\textbf{Proof of the reduction.} Before division by \(S_n\), the two
terms in the first coordinate have rates at most \(x_*+A_1\) and exactly
\(K-x_*+A_2\). Their difference is \(\gamma>0\), by (5.3), so the
Nesterenko term dominates even if the signs oppose. This proves the
first line of (6.3), including its lower bound. In the linear-form
value, the corresponding limiting rates are equal because of (6.1).
Taking \(x_n<x_*\) makes the Nesterenko term exponentially dominant by
\(2(x_*-x_n)n\), which proves the second, two-sided line of (6.3) and
also nonvanishing. Subtract (4.3) in both cases and apply (1.3).
\(\square\)

This reduction shows why a one-vector box argument, without a lower
height bound, does not prove (6.5). The unconditional argument below
obtains the height bound from two vectors instead.

\hypertarget{a-determinant-selection-theorem-requiring-no-annular-coordinate}{%
\subsubsection{A determinant selection theorem requiring no annular
coordinate}\label{a-determinant-selection-theorem-requiring-no-annular-coordinate}}

As an intermediate step, the following balanced-minima condition
isolates how two independent short vectors and the crossed determinant
produce the required lower height.

Let

\[
 \mathcal B_n=\{(s,t)\in\mathbb R^2:
 |s|\le e^{x_*n},\ |t|\le e^{(K-x_*)n}\},
 \tag{6.6}
\]

and let \(\lambda_{1,n}\le\lambda_{2,n}\) be the successive minima of
\(L_n\) with respect to \(\mathcal B_n\). Consider the two assertions

\[
 \log d_n=Kn+o(n),\qquad \log\lambda_{1,n}=o(n).
 \tag{BM}
\]

The first says that the extra common divisor \(g_n\) in (4.5) is
subexponential. The second excludes an exponentially exceptional first
minimum; its upper-bound half is automatic from Minkowski, while its
lower bound is new arithmetic information about the HNF residue (7.2).

\begin{quote}
\textbf{Balanced-minima selection theorem.} Assume (BM). Then there are
vectors \(c_n\in L_n\) for which the integer pairs in (4.7) satisfy \[
\log|q_n(c_n)|=A_{\rm eff}n+o(n),\qquad
0<|q_n(c_n)G-p_n(c_n)|\le e^{E_{\rm eff}n+o(n)}.       \tag{6.7}
\] Consequently \[
\delta(q,p)\ge\delta_*=0.5819834736933570952\ldots .  \tag{6.8}
\]
\end{quote}

\textbf{Proof.} The area of \(\mathcal B_n\) is \(4e^{Kn}\). Minkowski's
second theorem and (BM) imply

\[
 \log\lambda_{2,n}=o(n).
\]

Thus there are independent \(c_n^{(1)},c_n^{(2)}\in L_n\) such that, for
\(i=1,2\),

\[
 |c_{1,n}^{(i)}|\le e^{x_*n+o(n)},\qquad
 |c_{2,n}^{(i)}|\le e^{(K-x_*)n+o(n)}.                \tag{6.9}
\]

Their determinant is a nonzero multiple of \(d_n\), while (6.9) bounds
it above by \(e^{Kn+o(n)}\). Hence

\[
 \log|\det(c_n^{(1)},c_n^{(2)})|=Kn+o(n).             \tag{6.10}
\]

Write \(q_i=q_n(c_n^{(i)})\), \(p_i=p_n(c_n^{(i)})\), and
\(\ell_i=q_iG-p_i\). Equations (3.3), (4.3), and (6.9), together with
the strict coefficient gap (5.3), give

\[
 |q_i|\le e^{A_{\rm eff}n+o(n)},\qquad
 |\ell_i|\le e^{E_{\rm eff}n+o(n)}.                  \tag{6.11}
\]

On the other hand, (5.4) and (6.10) give the exact determinant rate

\[
 \begin{aligned}
 \log|q_1p_2-q_2p_1|
 &=\log\left|
 \frac{\det(c_n^{(1)},c_n^{(2)})\mathcal D_n}{S_n^2}
 \right|\\
 &=(A_2+E_1-K)n+o(n)\\
 &=(A_{\rm eff}+E_{\rm eff})n+o(n).
 \end{aligned}                                      \tag{6.12}
\]

Since

\[
 q_1p_2-q_2p_1=q_2\ell_1-q_1\ell_2,
\]

(6.11)--(6.12) force

\[
 \max(|q_1|,|q_2|)\ge e^{A_{\rm eff}n-o(n)}.         \tag{6.13}
\]

Choose an index \(i\) realizing (6.13). If \(\ell_i\ne0\), take
\(c_n=c_n^{(i)}\). If \(\ell_i=0\), independence of the two coefficient
vectors implies that the other linear form is nonzero; choose a sign
\(\varepsilon\in\{-1,1\}\) for which

\[
 |q_i+\varepsilon q_{3-i}|\ge |q_i|
\]

and take \(c_n=c_n^{(i)}+\varepsilon c_n^{(3-i)}\). This changes every
upper bound only by a constant factor and makes the linear form nonzero.
Together with (6.11) and (6.13), this proves (6.7), including the lower
height bound. Formula (1.1) now gives (6.8). \(\square\)

Unlike the one-vector Minkowski argument, this lemma survives accidental
cancellation in either source coordinate. The next theorem removes (BM)
entirely by treating an exponentially short first minimum as a useful
reduction direction.

\hypertarget{unconditional-reduction-in-the-linear-form-coordinate}{%
\subsubsection{Unconditional reduction in the linear-form
coordinate}\label{unconditional-reduction-in-the-linear-form-coordinate}}

The preceding conditional formulation is useful for orientation but is
not needed. When the first minimum is exponentially short, one can
reduce the second vector modulo the first in the \textbf{linear-form
value}, and the crossed determinant then forces the required height.
This closes the selection step for every possible shape of (L\_n).

Put

\[
 s_n=\frac{\log S_n}{n},\qquad
 \kappa_n=\frac{\log d_n}{n},\qquad
 x_n=\frac{\kappa_n+E_2-E_1}{2},                    \tag{6.14}
\]

and

\[
 H_n=\kappa_n-x_n+A_2-s_n,qquad
 F_n=x_n+E_1-s_n.                                   \tag{6.15}
\]

Thus \(s_n\to K=6\), \(0\le\kappa_n\le s_n\), and

\[
 H_n-F_n=A_2-E_2,qquad
 H_n+F_n=\kappa_n+A_2+E_1-2s_n.                    \tag{6.16}
\]

Both \(H_n\) and \(F_n\) are bounded away from zero. Since
\(\kappa_n\le s_n\),

\[
 \limsup_{n\to\infty}H_n\le A_{\rm eff},\qquad
 \limsup_{n\to\infty}F_n\le E_{\rm eff}.           \tag{6.17}
\]

Let

\[
 \mathcal C_n=\{(u,v)\in\mathbb R^2:
 |u|\le e^{x_nn},\ |v|\le e^{(\kappa_n-x_n)n}\}.
 \tag{6.18}
\]

Its area is \(4d_n\). If \(\mu_{1,n}\le\mu_{2,n}\) are the successive
minima of \(L_n\) with respect to \(\mathcal C_n\), Minkowski's second
theorem gives

\[
 \mu_{1,n}\mu_{2,n}\le1.                           \tag{6.19}
\]

Choose independent vectors \(b_n^{(1)},b_n^{(2)}\in L_n\) at these
successive minima. Their determinant is a nonzero multiple of \(d_n\),
whereas (6.19) bounds its absolute value by \(2d_n\). Hence

\[
 \log|\det(b_n^{(1)},b_n^{(2)})|=\log d_n+O(1).      \tag{6.20}
\]

Write

\[
 q_i=q_n(b_n^{(i)}),\quad p_i=p_n(b_n^{(i)}),
 \quad \ell_i=q_iG-p_i.                              \tag{6.21}
\]

\begin{quote}
\textbf{Theorem 6.2 (unconditional two-minimum selection).} For every
sufficiently large \(n\) one can select \(c_n\in L_n\) so that, on
writing \(q_n=q_n(c_n)\), \(p_n=p_n(c_n)\), and \(\ell_n=q_nG-p_n\),

\[
q_n\ne0,\qquad \ell_n\ne0,\qquad
\log|q_n|\ge H_nn-o(n),                            \tag{6.22}
\]

and

\[
\frac{\log|\ell_n|}{\log|q_n|}
\le \frac{F_n}{H_n}+o(1).                         \tag{6.23}
\]

Consequently \(|q_n|\to\infty\) and \[
\liminf_{n\to\infty}
\frac{-\log|G-p_n/q_n|}{\log|q_n|}
\ge\frac{A_2-E_2}{A_{\rm eff}}.                   \tag{6.24}
\]
\end{quote}

\textbf{Proof.} Choose \(\epsilon_n>0\), \(\epsilon_n\to0\), dominating
all normalized remainders in (3.3), (4.3), (5.4), and (6.20). Enlarging
it if necessary, assume also that \(n\sqrt{\epsilon_n}\to\infty\). Since
\(\mu_{1,n}\le1\) by (6.19), put

\[
 r_n=-\frac{\log\mu_{1,n}}n\ge0.
\]

By (6.19), the two vectors have the respective bounds obtained from
\(\mathcal C_n\) by common factors \(e^{-r_nn}\) and \(e^{r_nn}\).
Moreover, \(r_n=O(1)\): every nonzero integer vector has a nonzero
integral coordinate, while \(x_n\) and \(\kappa_n-x_n\) are bounded. The
definition of \(x_n\) balances the value terms, while the Nesterenko
coefficient term exceeds the Zudilin coefficient term by the fixed gap
\(\gamma\) in (5.3). Therefore

\[
 \begin{array}{ll}
 |q_1|\le e^{(H_n-r_n+O(\epsilon_n))n},&
 |\ell_1|\le e^{(F_n-r_n+O(\epsilon_n))n},\\
 |q_2|\le e^{(H_n+r_n+O(\epsilon_n))n},&
 |\ell_2|\le e^{(F_n+r_n+O(\epsilon_n))n}.
 \end{array}                                         \tag{6.25}
\]

The exact determinant identity is

\[
 \begin{aligned}
 \Delta_n'&:=q_1p_2-q_2p_1=q_2\ell_1-q_1\ell_2\\
 &=\frac{\det(b_n^{(1)},b_n^{(2)})\mathcal D_n}{S_n^2}.
 \end{aligned}                                       \tag{6.26}
\]

Equations (5.4), (6.16), and (6.20) imply

\[
 \log|\Delta_n'|=(H_n+F_n)n+O(\epsilon_nn).          \tag{6.27}
\]

If \(r_n\le\sqrt{\epsilon_n}\), (6.25)--(6.27) force

\[
 \max(|q_1|,|q_2|)\ge e^{(H_n-O(\sqrt{\epsilon_n}))n}.
\]

Choose the vector with the larger first coordinate. If its linear form
is zero, add or subtract the other vector, choosing the sign so that the
first coordinates do not cancel. The two independent vectors cannot both
have zero linear form. This gives a nonzero form and

\[
 \log|q_n(c_n)|=H_nn+o(n),\qquad
 \log|q_n(c_n)G-p_n(c_n)|\le F_nn+o(n).              \tag{6.28}
\]

Suppose \(r_n>\sqrt{\epsilon_n}\) and first that \(\ell_1\ne0\). Choose
an integer \(m_n\) satisfying

\[
 0<|\ell_2-m_n\ell_1|\le\frac32|\ell_1|;             \tag{6.29}
\]

a nearest integer to \(\ell_2/\ell_1\) works unless its remainder is
zero, in which case an adjacent integer works. Set

\[
 c_n=b_n^{(2)}-m_nb_n^{(1)},\quad
 q'=q_2-m_nq_1,\quad \ell'=\ell_2-m_n\ell_1.
\]

From (6.26),

\[
 q'=\frac{\Delta_n'}{\ell_1}
     +q_1\frac{\ell'}{\ell_1}.                       \tag{6.30}
\]

The first term has size at least \(e^{(H_n+r_n-O(\epsilon_n))n}\), while
the second is at most \(e^{(H_n-r_n+O(\epsilon_n))n}\). Since
\(r_n>\sqrt{\epsilon_n}\) and \(n\sqrt{\epsilon_n}\to\infty\), the first
term dominates by a factor tending to infinity. If
\(a_n=n^{-1}\log|\ell_1|\), then

\[
 \frac{-\log|G-(p_2-m_np_1)/(q_2-m_nq_1)|}
 {\log|q_2-m_nq_1|}
 \ge
 1-\frac{F_n-r_n+O(\epsilon_n)}
          {H_n+r_n-O(\epsilon_n)}.                  \tag{6.31}
\]

Indeed, (6.30) gives \(n^{-1}\log|q'|=H_n+F_n-a_n+O(\epsilon_n)\), and
the right side is worst when \(a_n=F_n-r_n+O(\epsilon_n)\); a smaller
\(\ell_1\) only improves the quotient. It also gives
\(\log|q'|\ge(H_n+r_n-O(\epsilon_n))n\). Finally,

\[
 \frac{F_n-r_n}{H_n+r_n}\le\frac{F_n}{H_n}
\]

because \(H_n,F_n>0\) and \(r_n\ge0\). Thus this branch satisfies both
(6.22) and (6.23).

If \(\ell_1=0\), then (6.26) implies \(q_1\ne0\) and \(\ell_2\ne0\).
This case also admits a completely fixed shear. For all sufficiently
large \(n\), \(5<s_n<7\); the closed forms (3.2) give \(1<E_2-E_1<2\),
\(13<E_1<14\), and \(29<A_2<30\). Consequently

\[
 \begin{aligned}
 \tfrac12&<x_n<\tfrac92,& -1&<\kappa_n-x_n<3,&
 r_n&<\tfrac92,\\
 H_n&<28,& F_n&<\tfrac{27}{2},&
 \frac{F_n}{H_n}&>\frac{13}{56}.
 \end{aligned}                                      \tag{6.32a}
\]

Here the bound for \(r_n\) follows because a nonzero integral vector has
a coordinate of absolute value at least one. Equation (6.25) now gives
\(\log^+|q_2|<34n\) and \(\log|\ell_2|<19n\) eventually. Take the fixed
shear

\[
 c_n=b_n^{(2)}+\lceil e^{100n}\rceil b_n^{(1)}.
\]

Since the nonzero integer \(q_1\) has absolute value at least one, this
gives \(\log|q_n(c_n)|\ge100n-o(n)\) and leaves the nonzero form equal
to \(\ell_2\). Thus its relative value ratio is at most
\(19/100+o(1)<13/56<F_n/H_n\). Notice that this exceptional case implies
\(G=p_1/q_1\in\mathbb Q\); it must be covered because the irrationality
of \(G\) is not known.

All three cases give the height bound (6.22), the relative value bound
(6.23), and hence

\[
 \frac{-\log|G-p_n(c_n)/q_n(c_n)|}{\log|q_n(c_n)|}
 \ge1-\frac{F_n}{H_n}-o(1)
 =\frac{A_2-E_2}{H_n}-o(1).                         \tag{6.32}
\]

Because \(H_n\) is bounded away from zero, the selected heights tend to
infinity. Equation (6.17) now proves (6.24). \(\square\)

The selection can be made canonical in classical exact-real arithmetic.
To make it a literal sequence rather than an existence choice, order
\(\mathbb Z^2\) by \((|u|+|v|,u,v)\), use the first vector attaining
each successive minimum (subject to independence for the second), break
all sign and nearest-integer ties by the same order, and apply the three
cases in the proof. The zero-form branch is retained explicitly because
deciding it algorithmically would amount to resolving a rationality
question about \(G\).

With

\[
 x_* =\frac{K+E_2-E_1}{2}
      =3.6467782381824059906\ldots,
\]

we have

\[
 \begin{aligned}
 A_{\rm eff}&=A_2-x_*=25.707995623896753637\ldots,\\
 E_{\rm eff}&=E_2-x_*=10.746367029007697991\ldots.
 \end{aligned}                                      \tag{6.33}
\]

Since \(A_2-E_2=A_{\rm eff}-E_{\rm eff}\), Theorem 6.2 and the
elementary identity \(\liminf(1-z_n)=1-\limsup z_n\) prove

\[
 \boxed{\delta(q,p)\ge
 1-\frac{E_{\rm eff}}{A_{\rm eff}}
 =0.5819834736933570952\ldots.}                     \tag{6.34}
\]

Explicitly, the final liminf/limsup passage is

\[
 \begin{aligned}
 \delta(q,p)
 &=1-\limsup_{n\to\infty}
       \frac{\log|q_nG-p_n|}{\log|q_n|}\\
 &\ge1-\limsup_{n\to\infty}\frac{F_n}{H_n}
 \ge1-\frac{E_{\rm eff}}{A_{\rm eff}}.
 \end{aligned}                                      \tag{6.35}
\]

The last inequality does not come from comparing separate limsups. It
uses the exact relation

\[
 \frac{F_n}{H_n}=1-\frac{A_2-E_2}{H_n},
\]

which is increasing as a function of \(H_n>0\), together with
\(\limsup H_n\le A_{\rm eff}\) from (6.17).

This conclusion is unconditional. No annular-coordinate assertion,
common gcd estimate, balanced-minimum hypothesis, or numerical
certificate is used.

\hypertarget{hnf-interpretation-of-the-superseded-one-vector-approach}{%
\subsection{7. HNF interpretation of the superseded one-vector
approach}\label{hnf-interpretation-of-the-superseded-one-vector-approach}}

For comparison, the earlier one-vector obstruction can be stated without
geometry-of-numbers jargon. Divide \(Y_n,U_n,S_n\) by \(g_n\), and write

\[
 Y'_n=Y_n/g_n,\quad U'_n=U_n/g_n,\quad d_n=S_n/g_n.
\]

Put

\[
 v_n^{\rm H}=\gcd(d_n,Y'_n),\qquad
 h_n^{\rm H}=d_n/v_n^{\rm H}.
 \tag{7.1}
\]

There is a unique \(a_n\), \(0\le a_n<h_n^{\rm H}\), satisfying

\[
 (Y'_n/v_n^{\rm H})a_n+U'_n\equiv0
       \pmod {h_n^{\rm H}}.
 \tag{7.2}
\]

Then

\[
 L_n=\mathbb Z(h_n^{\rm H},0)+
     \mathbb Z(a_n,v_n^{\rm H}),qquad
 h_n^{\rm H}v_n^{\rm H}=d_n.                       \tag{7.3}
\]

Consequently (AS) is equivalent to finding integers \(k_n,m_n\) with

\[
 c_{2,n}=k_nv_n^{\rm H},\qquad
 c_{1,n}=k_na_n-m_nh_n^{\rm H},                    \tag{7.4}
\]

such that (6.2) holds. This is an annular modular-approximation theorem
for the \textbf{specific} residues \(a_n/h_n^{\rm H}\).

The first minimum in (BM) has the equally explicit formula

\[
 \lambda_{1,n}=\min_{(k,m)\ne(0,0)}
 \max\left\{
   |ka_n-mh_n^{\rm H}|e^{-x_*n},
   |k|v_n^{\rm H}e^{-(K-x_*)n}
 \right\}.                                           \tag{7.5}
\]

Thus the lower half of the intermediate condition (BM) says that the
minimum in (7.5) is not \(e^{-\eta n}\) for any fixed \(\eta>0\),
eventually. It is a precise modular-approximation assertion, not a
consequence of the lattice index.

The index estimate (4.6) alone cannot prove it. For example,

\[
 L(S)=\{(c_1,c_2)\in\mathbb Z^2:c_2\equiv0\pmod S\}
\]

has index \(S\), but contains no vector with \(0<|c_2|<S\). Hence any
valid proof of (AS) must use additional arithmetic of \(Y_n,U_n\), or a
genuinely stronger selection theorem whose conclusion allows a
compensating alternative and still proves a lower height bound. Neither
ordinary Minkowski nor the existing Lean \texttt{ExponentPair} predicate
does this.

Row independence alone does not repair the defect. For example, the two
integer rows \((M,0)\) and \((M+1,1)\) have determinant \(M\ne0\), while
the short combination \((1,-1)\) has first coordinate \(-1\) although
its two summands have size comparable with \(M\). Thus a proposed
lower-height lemma must also survive nearly proportional rows; nonzero
determinant plus coefficient upper bounds is not enough.

\hypertarget{like-for-like-benchmarks}{%
\subsection{8. Like-for-like
benchmarks}\label{like-for-like-benchmarks}}

The decimal comparisons below can be certified without floating-point
assumptions. For completeness, bracket each square root by rationals
whose squares bracket the radicand, rescale every positive logarithm
argument to \(y\in[1,2)\), and use

\[
 \log y=2\sum_{j=0}^{N}\frac{z^{2j+1}}{2j+1}+R_N,qquad
 z=\frac{y-1}{y+1},qquad
 0<R_N<\frac{2z^{2N+3}}{(2N+3)(1-z^2)}.             \tag{8.1}
\]

For example, the exact integer square comparisons give

\[
 \begin{aligned}
 \frac{2236067977499789}{10^{15}}&<\sqrt5<
 \frac{2236067977499790}{10^{15}},\\
 \frac{7549834435270749}{10^{15}}&<\sqrt{57}<
 \frac{7549834435270750}{10^{15}}.
 \end{aligned}
\]

Taking \(N=40\) and rational arithmetic gives the outward-rounded
intervals

\[
\begin{array}{rcl}
27.53594&<&A_1<27.53595,\\
13.09958&<&E_1<13.09960,\\
29.35477&<&A_2<29.35478,\\
14.39314&<&E_2<14.39315,\\
0.58198347&<&1-E_{\rm eff}/A_{\rm eff}<0.58198348,\\
0.52427310&<&1-E_1/A_1<0.52427312,\\
0.50968297&<&1-E_2/A_2<0.50968299.
\end{array}                                          \tag{8.2}
\]

Thus the displayed comparisons have a purely rational interval
certificate; the many-digit decimals are not being used as proof.

In the same definition (1.1):

\begin{itemize}
\tightlist
\item
  Zudilin's asymptotic square-cleared row has \[
  1-E_1/A_1=0.5242731097791163788\ldots.
  \]
\item
  The unconditional conservative Nesterenko row (2.14) has quality \[
  1-E_2/A_2=0.5096829791700988894\ldots.
  \]
\item
  Under the unproved integrality of Nesterenko's sharper numerator, his
  sharper normalization has limiting log-ratio
  \(0.4455574797507435\ldots\), hence limiting construction quality
  \(0.5544425202492564\ldots\). The paper states the safe conditional
  inequality with exponent \(11/20=0.55\), not \(0.5545\).
\end{itemize}

The unconditional bound (6.34) exceeds all three comparison values. The
comparison is like for like: every number is a construction quality in
the sense of (1.1), rather than a classical irrationality exponent for
\(G\).

\hypertarget{primary-sources}{%
\subsection{9. Primary sources}\label{primary-sources}}

\begin{enumerate}
\def\labelenumi{\arabic{enumi}.}
\tightlist
\item
  W. Zudilin, \emph{An Apéry-like difference equation for Catalan's
  constant}, \textbf{Electronic Journal of Combinatorics} 10 (2003),
  Research Paper 14;
  \href{https://arxiv.org/abs/math/0201024}{arXiv:math/0201024}.
\item
  W. Zudilin, \emph{A few remarks on linear forms involving Catalan's
  constant}, \textbf{Chebyshevskii Sbornik} 3:2 (2002), 60--70;
  \href{https://arxiv.org/abs/math/0210423}{arXiv:math/0210423}.
\item
  Yu. V. Nesterenko, \emph{On Catalan's constant}, \textbf{Proceedings
  of the Steklov Institute of Mathematics} 292 (2016), 153--170,
  \href{https://doi.org/10.1134/S0081543816010107}{doi:10.1134/S0081543816010107}.
\end{enumerate}

Exact page and equation references, including the two printed
normalization/sign issues in the third source, are recorded in
\texttt{PRIMARY\_SOURCE\_AUDIT.md}.

\end{document}
